\documentclass[11pt,a4paper]{article}
\usepackage[utf8]{inputenc}
\usepackage[T1]{fontenc}
\usepackage[italian]{babel}
\usepackage[a4paper,top=1.8cm,bottom=1.8cm,left=2cm,right=2cm,headheight=15pt]{geometry}
\usepackage{amsmath,amssymb}
\usepackage{graphicx}
\usepackage{tikz}
\usepackage{pgffor}
\usepackage{enumitem}
\usepackage{fancyhdr}
\usepackage{lastpage}
\usepackage{booktabs}
\usepackage{longtable}
\usepackage{siunitx}
\usepackage{xcolor}
\usepackage{needspace}
\usepackage{array}
\usepackage{textcomp}
\usepackage{pdflscape}
\usepackage[hidelinks]{hyperref}
\sisetup{output-decimal-marker={,},group-separator={.},per-mode=symbol}
\setlength{\parindent}{0pt}
\setlength{\parskip}{2pt}
\setlist[enumerate]{itemsep=6pt,topsep=4pt,leftmargin=8mm}
\setlist[itemize]{itemsep=2pt,topsep=3pt,leftmargin=6mm}
\raggedbottom

% Comandi personalizzati: modifica qui la grafica senza toccare il codice.
\pagestyle{fancy}
\fancyhf{}
\renewcommand{\headrulewidth}{0.4pt}
\fancyfoot[C]{\small Pag. \thepage\ di \pageref{LastPage}}

% Punteggio a destra della consegna.
\newcommand{\punti}[1]{\hfill\textbf{[#1~pt]}}

% Riga puntinata per lo svolgimento.
\newcommand{\riga}{\noindent\makebox[\linewidth]{\dotfill}\par\vspace{4mm}}

% N righe puntinate (se N = 0 non stampa nulla: il chiamante evita il caso).
\newcommand{\rispostarighe}[1]{\par\vspace{1mm}\foreach \i in {1,...,#1}{\riga}}

% Riquadro per la risposta.
\newcommand{\riquadro}{\par\noindent\framebox[\linewidth]{\rule{0pt}{2.6cm}}\par}

% Griglia quadrettata per grafici e disegni.
\newcommand{\griglia}{\par\noindent
  \begin{tikzpicture}[x=0.55cm,y=0.55cm]
    \draw[step=1,lightgray,very thin] (0,0) grid (16,10);
    \draw[black!60,thick] (0,0) rectangle (16,10);
  \end{tikzpicture}\par}

% Piano cartesiano con assi e scala.
\newcommand{\grigliacartesiana}{\par\noindent
  \begin{tikzpicture}[x=0.55cm,y=0.55cm]
    \draw[step=1,lightgray,very thin] (0,0) grid (16,10);
    \draw[->,thick] (0,5) -- (16,5) node[right] {\small $x$};
    \draw[->,thick] (8,0) -- (8,10) node[above] {\small $y$};
    \foreach \x in {2,4,6,10,12,14} \draw (\x,4.8) -- (\x,5.2);
    \foreach \y in {1,3,7,9} \draw (7.8,\y) -- (8.2,\y);
  \end{tikzpicture}\par}

% Intestazione della verifica: istituto, indirizzo, materia, classe, data, tempo, punti, soglia.
\newcommand{\intestazione}[8]{%
  \begin{center}
    {\large\bfseries #1}\\[1mm]
    {\small #2\ --\ #3}\\[3mm]
    {\Large\bfseries #4}\\[3mm]
    \begin{tabular}{@{}p{0.48\textwidth}p{0.48\textwidth}@{}}
      Classe: #5 & Data: #6 \\
      Tempo: #7 & Punteggio: #8 \\
    \end{tabular}
  \end{center}
}

% Titolo di sezione con spazio garantito.
\newcommand{\titolosezione}[1]{\par\needspace{4\baselineskip}\section*{#1}\vspace{-1mm}}

% Blocco soluzione.
\newcommand{\soluzione}[1]{\par\noindent #1\par}

% Casella per gli errori comuni (soluzioni).
\newcommand{\erroricomuni}[1]{%
  \par\smallskip
  \noindent\fcolorbox{black!25}{black!3}{%
    \begin{minipage}{\dimexpr\linewidth-2\fboxsep-2\fboxrule\relax}
      \small\textbf{Errori comuni attesi:} #1
    \end{minipage}}\par\smallskip
}


\begin{document}

\fancyhead[L]{\small Liceo Statale di Esempio}
\fancyhead[R]{\small Matematica\ --\ 12/11/2026}

\intestazione{Liceo Statale di Esempio}{Liceo Classico}{Matematica}{Verifica dimostrativa --- Equazioni e geometria (1\textsuperscript{a} Classico)}{1\textsuperscript{a} A}{12/11/2026}{60}{100}




\noindent Nome e cognome: \dotfill \hfill Classe: 1\textsuperscript{a} A \hfill Voto: \dotfill
\par\vspace{1mm}
\noindent\rule{\linewidth}{0.4pt}
\par\vspace{1mm}

\noindent\small Non è consentito l'uso della calcolatrice. Scrivi lo svolgimento negli spazi indicati.\par\vspace{2mm}



\titolosezione{Quesiti a risposta chiusa}

\noindent\small Una sola risposta corretta.\par\vspace{1mm}

\begin{enumerate}[label=\arabic*),start=1]
\needspace{5\baselineskip}
\item Trova il valore di $x$ che rende vera l'uguaglianza $3x - 8 = 10$. Scegli la risposta corretta. \punti{4,5}
\begin{enumerate}[label=\Alph*),itemsep=1pt,topsep=2pt,leftmargin=7mm]

  \item $x = 6$

  \item $x = -6$

  \item $x = 5$

  \item $x = 7$

\end{enumerate}

\needspace{5\baselineskip}
\item Scomponi il trinomio $x^2 + 2x + 1$. Scegli la risposta corretta. \punti{4,5}
\begin{enumerate}[label=\Alph*),itemsep=1pt,topsep=2pt,leftmargin=7mm]

  \item $(x + -2)(x + 1)$

  \item $(x + 1)(x + 1)$

  \item $(x - 1)(x - 1)$

  \item $(x + 1)(x - 1)$

\end{enumerate}

\end{enumerate}

\titolosezione{Quesiti a risposta breve}

\begin{enumerate}[label=\arabic*),resume]
\needspace{5\baselineskip}
\item Un articolo costa 190 euro. Viene applicato uno sconto del 50\%. Calcola lo sconto e il prezzo finale. Scrivi il risultato e una breve giustificazione. Giustifica il procedimento seguito. \punti{6,75}

\rispostarighe{3}


\needspace{5\baselineskip}
\item Risolvi l'equazione $2x + 1 = 5$. Scrivi il risultato e una breve giustificazione. Giustifica il procedimento seguito. \punti{6,75}

\rispostarighe{3}


\end{enumerate}

\titolosezione{Esercizi}

\begin{enumerate}[label=\arabic*),resume]
\needspace{6\baselineskip}
\item Data la funzione $y = 3x - 4$, calcola il valore di $y$ per $x = 2$. Giustifica il procedimento seguito. \punti{13,75}


\rispostarighe{5}


\needspace{8\baselineskip}
\item Un sacchetto contiene 3 palline rosse e 2 palline blu. Estraendo una pallina a caso, qual è la probabilità che sia rossa? Giustifica il procedimento seguito. \punti{13,75}


\rispostarighe{10}


\end{enumerate}

\titolosezione{Problemi e dimostrazioni}

\begin{enumerate}[label=\arabic*),resume]
\needspace{12\baselineskip}
\item In un gioco, il punteggio finale si ottiene moltiplicando per 5 il numero di risposte corrette e togliendo 9 punti. Il punteggio finale è 16. Quante risposte corrette sono state date? Spiega perché il procedimento scelto è adeguato. \punti{31,75}


\rispostarighe{17}


\needspace{11\baselineskip}
\item Spiega perché l'area di un rettangolo si calcola moltiplicando base e altezza, usando come esempio un rettangolo di $4\,\mathrm{cm}$ per $9\,\mathrm{cm}$. \punti{18,25}

\rispostarighe{16}


\end{enumerate}


\end{document}
